\section*{Visão Geral da Estatística e Análise de Dados}

\subsection*{O que é Estatística}

A estatística é a ciência matemática -- e parte da metodologia científica -- que busca
extrair, reduzir, analisar e modelar um conjunto de dados.
Seu principal objetivo é encontrar padrões em dados e possibilitar a tomada de conclusões
e decisões bem informadas sob incerteza.
A Teoria das Probabilidades, que estuda os eventos e as variáveis aleatórias, é
a principal base matemática dos conhecimentos estatísticos, conferindo-lhes rigor e
confiabilidade.

Para encontrar tais padrões num conjunto de dados, a estatística faz o uso de modelos
estatísticos.
Esses modelos tentam descrever o dado em duas partes: a parte suave, onde se encontram
os padrões, e a parte ruidosa, que corresponde ao fator aleatório inerente a fonte
de extração do conjunto.
A formulação dos modelos estatísticos ainda pode ser realizada visando-se dois fins,
sendo o primeiro o de predizer o comportamento de dados não observados, enquanto que o
segundo consiste na inferência, isto é, tomar conclusões sobre os parâmetros encontrados.

A inferência estatística possui um lugar de centralidade no ramo estatístico, pois é sua
abordagem que permite generalizar as conclusões sobre os parâmetros dos modelos ajustados
sobre uma amostra para uma população, ao mesmo tempo que quantifica a incerteza dessas
conclusões.
Há duas grandes escolas de inferência estatística, que se diferenciam quanto a
interpretação do que é uma probabilidade e, consequentemente, como tratar a incerteza.
A primeira é a clássica ou frequentista, que entende probabilidades como frequências
relativas obtidas na observação a longo prazo, enquanto os parâmetros são considerados
fixos e desconhecidos, com a incerteza vindo da distribuição dos valores dados pelos
estimadores.
A segunda é a bayesiana, que interpreta probabilidade como o grau de crença num evento
ocorrer e vê os parâmetros como variáveis aleatórias, possuindo estes uma distribuição
própria atualizada pelas observações via o famoso Teorema de Bayes.

\subsection*{Uma Breve História da Estatística}

A palavra estatística deriva dos termos neolatim \textit{statisticum collegium} e italiano
\textit{statista}, que significam Colégio de Estado e Estadista respectivamente.
A palavra alemã \textit{Statistik} foi usada modernamente por Gottfried Achenwall
(1719--1772) em meados do século XVIII para designar a coleta de dados políticos,
econômicos e sociais de um país, cujo equivalente em inglês era \textit{political
arithmetic}.
Posteriormente, através de Sir John Sinclair (1754--1835) e seus trabalhos no
\textit{Statistical Accounts of Scotland}, a palavra ganhou uma conotação mais ampla sobre
coleta e classificação de dados.

Como vistos pelas suas raízes etimológicas, a estatística surgiu enquanto a disciplina
que se preocupava com a coleta, tratamento, classificação e análise de dados sobre
características político-sócio-econômicas de um país a fim de apoiar a formulação de
políticas pelos seus governos.
Dado isso, a prática estatística apareceu, provavelmente, há milhares de anos atrás,
quando foram realizados os primeiros censos e contagens de recursos pelos povos.
No entanto, enquanto disciplina científica, a estatística desenvolve-se a partir
do século XIX, apoiada pelas ferramentas matemáticas estudadas em séculos anteriores,
especialmente por aquelas da Teoria das Probabilidades.

O estudo moderno e documentado de probabilidades inicia-se no século XVI com Girolamo
Cardano (1501--1576), um polimata italiano que, motivado pelo seu hábito de apostar, escreveu o livro
``\textit{Liber de ludo aleae}'' (O livro dos jogos de azar) entre 1550 e 1564.
Nesse trabalho, ele documenta a técnica de computar probabilidades como razões
de resultados favoráveis sobre o número total de resultados possíveis, desde que todas
as conclusões possíveis sejam equiprováveis.
Entretanto, o livro não foi devidamente publicado até 1663, quando cálculos sobre jogos
de azar simples já estavam amplamente disseminados pela Europa.

O estudo de jogos de azar e, consequentemente, de probabilidades foi também realizado ao
longo das correspondências entre os franceses Pierre de Fermat (1601--1665) e Blaise
Pascal (1623--1662) em meados do século XVII, principalmente em apoio aos jogos do
nobre Chevalier de Méré (1607--1684).
Suas contribuições consistiram no uso de técnicas combinatórias para resolver diversos
tipos jogos, essas que se tornaram amplamente difundidas e compuseram as bases da moderna
Teoria de Probabilidades.

A influência dos estudiosos franceses estendeu-se ao holandês Christiaan Huygens
(1629--1695), que sistematizou os conhecimentos desenvolvidos naquele considerado o primeiro tratado
sistemático sobre Probabilidades: ``\textit{Van Rekeningh in Spelen van Geluck}''
(Sobre o Cálculo em Jogos de Azar), introduzindo conceitos como a esperança matemática
e registrando a resolução do ``Problema dos Pontos''.
Os trabalhos de Huygens também encontraram aplicações nos cálculos de taxas de anuidades
financeiras na Holanda, sobre as quais também publicou um tratado em 1671.
O uso vasto e sofisticado de probabilidades para cálculos de riscos, taxas de empréstimos
e de seguros se deu mais futuramente em meados do século XVIII na Grã-Bretanha.

Em 1713, é publicada postumamente a célebre obra \textit{Ars conjectandi}, de Jakob
Bernoulli (1655--1705), reunindo e aprofundando ideias dos autores anteriormente citados
e dando a primeira formulação da Lei dos Grandes Números.
No ano de 1718, Abraham de Moivre (1667--1754) publica \textit{The Doctrine of Chances},
outro marco para a Teoria de Probabilidades, trazando uma definição formal de
independência estatística e apresentando a primeira formulação do Teorema Central do
Limite, a qual consistia no caso limite da distribuição binomial.
Já em 1763, é enviado, postumamente, um ensaio de Thomas Bayes (1701--1761) à \textit{Royal Society}
britânica que discutia o hoje conhecido Teorema de Bayes, que introduz o conceito de
probabilidade inversa e atualização de probabilidades à luz de novas evidências.

Na França, a probabilidade ainda encontrou aplicações nas chamadas ciências morais, no
sistema judiciário e em eleições de cargos.
Pierre de Laplace (1749--1827) e o Marquês de Condorcet (1743--1794), usaram as probabilidades para calcular as
chances de um veredito de um júri ou um grupo de eleitores tomarem a decisão ``correta''.
Em 1812, Laplace publica \textit{Théorie analytique des probabilités}, um clássico
tratado que expandia o estudo das probabilidades por meio das técnicas do Cálculo
Diferencial e Integral desenvolvidas no século anterior.
Uma das contribuições desse trabalho é a formulação e prova geral do Teorema Central do
Limite.

Paralelamente ao intenso desenvolvimento da Teoria de Probabilidades, a estatística moderna
começava a ganhar forma ao longo dos séculos XVII e XVIII através do nome de
\textbf{aritmética política}.
Dado o contexto do Iluminismo e o interesse em generalizar o comportamento social
por meio de leis universais, assim como a física fizera com a natureza, houve um
crescimento nos estudos sobre população, doenças e renda.
Exemplos desses estudos incluem os de John Graunt (1620--1674) e William Petty
(1623--1687), que publicaram, respectivamente, registros de mortalidade da população
londrina e da renda irlandesa, lançando as bases da estatística empírica.
O pastor prussiano Johann Peter Süssmilch (1707--1767) também publicou dados sobre proporções
de nascimentos e mortes, atrelados a ideia da manifestação da Divina Providência
na regularidade demográfica mostrada pelos dados.
Todavia, a aritmética política ainda carecia de métodos de coleta mais gerais sobre as
populações.

Em 1749, a Suécia torna-se a primeira a nação a realizar a coleta regular e contínua de
dados demográficos, primeiramente através de dados providos pela Igreja Luterana e depois
por meio de um instituto nacional em 1858, sendo motivada por razões militares e econômicas.
A constituição dos Estados Unidos inova ao trazer no seu texto a obrigatoriedade da
promoção de censos demográficos a cada 10 anos a partir de 1790, cuja motivação era de
atualizar o número de deputados para o Congresso e a distribuição dos impostos federais.
Muitas nações europeias também vieram a adotar a prática dos censos em meados do século
XIX.

Com as novas demandas criadas pelos censos e a maior oferta de dados para análise,
a estatística cresceu e floresceu ao longo do século XIX, sendo aplicada nos mais
diversos campos: dados territoriais, produção agrícola, renda, nascimentos, mortes,
criminalidade, suicídios, entre outros.
Um dos principais nomes desse movimento foi Adolphe Quetelet (1796--1874), que cunhou o conceito do
``Homem Médio'' para descrever centralidade e variações de características físicas e
morais, demonstrando também a regularidade estatística de fenômenos sociais quando
observados para uma grande população.
Quetelet também articulou o Primeiro Congresso Internacional de Estatística de 1853 em
Bruxelas, o qual presidiu e onde advogou pela uniformização, modernização e integração
da nova ciência estatística.

Concomitantemente ao desenvolvimento da aritmética política, outras ciências matemáticas
contribuiram para a formação da disciplina estatística.
Roger Cotes (1682--1716) publicou \textit{Opera Miscellanea} em 1722, inaugurando a
Teoria de Erros através da análise astronômica e o uso da média para combinar observações.
Thomas Simpson (1710--1761) aplica a análise probabilística aos erros de observação em 1755,
argumentando formalmente o porquê a média é mais confiável que um dado isolado.

Adrien-Marie Legendre (1752--1833) publica o Método dos Mínimos Quadrados em 1805,
utilizando-o na geodésia para estimar o comprimento do meridiano da Terra e determinar a
unidade do metro.
Carl Friedrich Gauss (1777--1855) desenvolve e aplica a técnica de forma independente para traçar a
órbita de Ceres em 1809, constatando também que ela encontrava o parâmetro ótimo sob a
hipótese dos erros seguirem a hoje denominada distribuição Normal ou Gaussiana.
Adiante, o Método dos Mínimos Quadrados tornou-se amplamente difundido para obter
estimativas de parâmetros de populações a partir de amostras em várias disciplinas
quantitativas.

Na Teoria Cinética dos Gases, a regularidade nas estatísticas sociais inspirou modelos
termodinâmicos desenvolvidos nos trabalhos de Rudolf Clausius (1822--1888), James Clerk
Maxwell (1831--1879) e Ludwig Boltzmann (1844--1906).
Nesses trabalhos, é inaugurada a moderna Mecânica Estatística, que usa modelos
probablísticos para explicar propriedades macroscópicas de um sistema a partir da
observação do comportamento geral de suas partes microscópicas.

A biologia também trouxe grandes contribuições à estatística na forma da biometria.
Sir Francis Galton (1822--1911) utilizou os conhecimentos estatísticos desenvolvidos até
meados do século XIX para analisar a diversidade biológica implicada pela Teoria da
Evolução de Charles Darwin (1809--1882).
Ele também foi um dos primeiros a desenvolver os métodos de regressões à média e
correlação, os quais aplicou à análise de hereditariedade.
Os trabalhos de Galton também influenciaram Karl Pearson (1857--1936), cujo nome batiza um dos
principais testes e medidas de associação em análise de dados, além de introduzir
o método dos momentos para estimação de parâmetros.
Pearson também funda o primeiro departamento universitário de estatística do mundo no
University College London em 1911, além da influente revista \textit{Biometrika}, que
durante muito tempo foi a principal difusora de novos métodos estatísticos e de
suas aplicações.

No início do século XX, observou-se a necessidade de métodos sistemáticos de planejamento
de experimentos e coleta de amostras que garantissem a eliminação de vieses e a
representatividade da população amostrada em várias variáveis.
Com esse fim, o estatístico polonês Jerzy Neyman (1894--1981) introduz o método de
amostragem estratificada em 1934.
O método ganha fama com George Gallup (1901--1984) em 1936, o qual previu a reeleição
de Franklin D. Roosevelt (1882--1945) à presidência dos Estados Unidos, em oposição a
previsão feita pela grande revista da época \textit{Literary Digest}, que apesar de
coletar um número maior de observações, o fez de forma a coletar apenas a opinião de
americanos de classe média alta, os quais tendiam a votar no candidato republicano
Alfred Landon (1887--1987).
A anedota mudou para sempre as noções do que seria uma amostragem aleatória e o design de
experimentos que necessitam de amostras representativas de populações.

Ainda na primeira metade e meados do século XX, o nome do biólogo e estatístico Ronald
Fisher (1890--1962) ganha destaque ao inaugurar a inferência estatística contemporânea.
Fisher expande o projeto e delineamento de experimentos, sistematiza o uso do valor-p para
significância estatística, cria o método ANOVA para estudos multipopulacionais
e introduz o método da Máxima Verossimilhança para estimação de parâmetros.
Fisher é um dos grandes nomes na sistematização de estudos quantitativos, contribuindo
para que a estatística não se prendesse somente a mera identificação de correlações, mas
indo também em direção às causas dos fenômenos, tendo contribuições na ecologia,
psicologia e pesquisas clínicas.

A Inferência Estatística ainda ganha diversas contribuições ao longo do século XX.
Em 1908, William Sealy Gosset (1876--1937), sob o pseudônimo de Student, publica, na
\textit{Biometrika}, a distribuição e o teste $t$, que permitiam inferência confiável a
partir de amostras pequenas, úteis à análise de dados na agricultura, indústria e
pesquisa clínica.
Jerzy Neyman e Egon Pearson (1895--1980), alternativamente à abordagem de Fisher, desenvolvem a
inferência a partir da formalização dos testes de hipótese, introduzindo conceitos
como erros do Tipo I e II, poder do teste e intervalos de confiança.

Na segunda metade do século XX, com o advento dos computadores, a estatística ganhou
novas fronteiras visto ao enorme poder de cálculo adquirido.
Nesse contexto, John Tukey (1915--2000) advoga por métodos mais flexíveis e visuais aliados aos
algoritmos computacionais para descobrir padrões ocultos e formular, num processo
iterativo, hipoteses sobre os dados, em oposição aos modelos estatísticos pré-concebidos
antes da análise.
Assim, Tukey inaugura a Análise Exploratória de Dados, a partir da qual também se
desenvolve a moderna Ciência de Dados do século XXI.

Atualmente a estatística é a disciplina central em qualquer outra área para lidar com
dados quantitativos e incertezas, enquanto que o profissional de estatística é a principal
referência para coleta, modelagem e análise de dados.
O avanço da tecnologia computacional viabilizou o desenvolvimento de técnicas e aplicações
estatísticas cada vez mais robustas e complexas para lidar com quantidades massivas de
dados e geração de conhecimento e compreensão em praticamente todas as ciências: físicas,
biológicas, sociais e humanas.
Muitas outras áreas também utilizam amplamente a estatística em suas práticas, como
engenharia, negócios, medicina, direito e atuária, por exemplo.

Visto seu desenvolvimento, conclui-se que a estatśitica permitiu a nós econtrar padrões e
regularidade onde, à primeira vista, só há desordem e ruído.
Ela representa o triunfo da humanidade sobre a incerteza acerca dos fatos mudanos.

\subsection*{Análise de Dados}

A Análise de Dados consiste na inspeção, limpeza, transformação e modelagem de dados
com o objetivo claro de obter informações úteis, informar conclusões ou apoiar a
tomada de decisão. Em resumo, ela é o processo de extração de dados brutos para
geração de conhecimento acionável, geralmente feito num ciclo iterativo. As fases desse
processo podem ser divididas em:

\begin{enumerate}
	\item Requisitos de Dados: a escolha de dados a serem extraídos dependem das
	      análises que se deseja realizar. Aqui define-se a \textbf{unidade experimental},
	      que pode consistir de um indivíduo ou uma população deles, e as \textbf{variáveis}
	      específicas dessa unidade, como idade e renda para uma população de pessoas.

	\item Coleta de Dados: consiste na tarefa de obter os dados de diferentes fontes,
	      como bancos de dados estruturados ou não-estruturados, sensores físicos, formulários,
	      uma interface de programação de aplicação (API), arquivos binários ou de textos, entre
	      muitas outras fontes.

	\item Processamento e Limpeza: inclui tarefas de organizar os dados em formatos
	      estruturados, identificar tipos, deduplicar, tratar valores faltantes, incorretos
	      e desviantes e transformar o formato de dados para outros mais convenientes.

	\item Análise Exploratória: compreende uma gama de técnicas para descrever e resumir
	      os dados obtidos à procura de informações que ajudem entender relações entre as
	      variáveis e guie às técnicas a serem aplicadas nas fases posteriores. Aqui
	      é comumente aplicadas técnicas da Estatística Descritiva e métodos gráficos.

	\item Modelagem: são usados algoritmos e modelos estatísticos para extrair padrões
	      latentes nos dados, muitas vezes não triviais de serem identificados. As
	      técnicas de aprendizado estatístico são especialmente úteis nessa etapa.

	\item Comunicação e Produto de Dados: a depender do contexto e das finalidades definidas
	      a priori, busca-se maneiras de comunicar esses dados aos atores interessados na
	      forma de visualizações e relatórios. Nos últimos anos, a disciplina de
	      \textbf{Narrativa de Dados} tem se destacado para realizar essa comunicação.
\end{enumerate}

A Análise de Dados, do ponto de vista da Estatística, poderia ser dividida em
\begin{enumerate}
	\item Estatística Descritiva: utilizar medidas de resumo numéricas para os dados
	      a fim de descrever padrões visíveis, geralmente respondendo perguntas do tipo
	      ``o que aconteceu?''.

	\item Análise Exploratória: semelhante a anterior, porém empregando métodos mais
	      flexíveis e visuais que permitam identificar padrões mais ocultos, correlações e
	      anomalias. Ocorre num processo iterativo de formulação de perguntas, exploração
	      e descobertas.

	\item Análise Confirmatória: aplica as técnicas da inferência estatística para
	      confirmar ou rejeitar hipóteses existentes das etapas anteriores
\end{enumerate}
